\chapter{INTRODUCTION}

% ============================================================================
% Section 1.1: Background of the Study
% ============================================================================
\section{Background of the Study}

Science education at the elementary level serves as a foundational stage in developing scientific literacy, conceptual reasoning, and problem-solving skills. Within the MATATAG Curriculum of the Department of Education (DepEd), Grade 6 Science introduces learners to fundamental concepts related to materials, including the description of changes of state for solids, liquids, and gases---specifically melting, evaporation, freezing, and condensation---using diagrams and flowcharts. These competencies aim to help learners understand observable phenomena through scientific representations and structured reasoning.

Despite the importance of these competencies, research has consistently shown that students experience difficulty in understanding the particulate nature of matter and the mechanisms underlying phase changes. Learners often memorize definitions of melting or evaporation without fully comprehending the particle-level interactions that explain these transformations. Misconceptions frequently arise regarding particle spacing, motion, and the role of heat transfer in altering states of matter (Adadan, 2013; Paik, Kim, Cho, \& Park, 2004). These challenges are compounded by the abstract nature of microscopic representations, which are not directly observable.

In many classroom settings, demonstrations of phase change are limited to water-based experiments, such as melting ice or boiling water. While these activities provide observable macroscopic changes, they often lack explicit integration with particle-level models, resulting in fragmented conceptual understanding. Furthermore, access to laboratory materials for demonstrating processes such as sublimation or deposition is limited at the elementary level, restricting exposure to diverse examples of state change.

\begin{figure}[htbp]
    \centering
    \begin{subfigure}[b]{0.32\textwidth}
        \centering
        \includegraphics[width=\textwidth]{PhaseChangeDiagram.jpg}
        \caption{States of matter and phase changes}
        \label{fig:phase-change-diagram}
    \end{subfigure}
    \hfill
    \begin{subfigure}[b]{0.32\textwidth}
        \centering
        \includegraphics[width=\textwidth]{KhanAcademyPhaseChanges.png}
        \caption{Phase changes (Khan Academy)}
        \label{fig:khan-academy-phases}
    \end{subfigure}
    \hfill
    \begin{subfigure}[b]{0.32\textwidth}
        \centering
        \includegraphics[width=\textwidth]{88GuruPhaseChanges.png}
        \caption{Particle arrangement in phases (88Guru)}
        \label{fig:88guru-phases}
    \end{subfigure}
    \caption{Examples of phase change diagrams used in science education}
    \label{fig:phase-changes-examples}
\end{figure}

Game-Based Learning (GBL) has emerged as a promising instructional approach that integrates interactive systems with curricular objectives. Educational games can provide dynamic visualizations, immediate feedback, and experiential learning opportunities that reduce cognitive overload and promote conceptual change (Plass, Homer, \& Kinzer, 2015). When designed with explicit learning outcomes and aligned assessments, digital games may enhance engagement and deepen understanding, particularly in abstract scientific domains.

This study introduces \textit{Wiz's Journey: Mastering the States of Matter}, a 2D educational platformer developed as a supplementary and enrichment learning tool aligned with the MATATAG Grade 6 Science competency on changes of state. The game integrates explicit terminology (e.g., ``Solid $\rightarrow$ Liquid: Melting'') while allowing learners to discover appropriate phase changes through environmental problem-solving. Central to its design is an interactive particle-lattice minigame---a core mechanic of the gameplay---in which learners manipulate bonds and particle spacing to simulate state transitions. By directly engaging with representations similar to those used in classroom diagrams, learners repeatedly encounter visual models of particle behavior in a contextualized, goal-driven environment.

Given the increasing emphasis on learner-centered instruction under the MATATAG framework, there is a need to explore innovative digital tools that complement formal instruction without replacing it. This study examines whether a platformer-based educational game can support conceptual understanding, engagement, and reinforcement of curriculum-aligned competencies related to phase changes of matter.


% ============================================================================
% Section 1.2: Objectives
% ============================================================================
\section{Objectives}

% Subsection 1.2.1: General Objective
\subsection{General Objective}

% TODO: Add general objective here

% Subsection 1.2.2: Specific Objectives
\subsection{Specific Objectives}

% TODO: Add specific objectives as enumerated list


% ============================================================================
% Section 1.3: Scope and Limitations of the Study
% ============================================================================
\section{Scope and Limitations of the Study}

% TODO: Add scope and limitations content here


% ============================================================================
% Section 1.4: Significance of the Study
% ============================================================================
\section{Significance of the Study}

% TODO: Add significance content here
